% answers_short.tex —— 2012 年·四川卷·文科 · 简短答案速查
% 由 tools/make_short_answers.py 自 answers.tex 的 \gksolhead/\gkanswer 提取，
% 并与 2012/2012_数学真题_四川文_答案要点.md 逐题比对后生成，勿手改（改源 MD 后重跑）。

\begin{gknotice}
\textbf{使用说明}：本节只给最终答案便于快速核对，完整推导见解析版。答案为编写者独立推导，\textbf{非官方标准答案}。
\end{gknotice}

\gkpart{选择题}
\begin{center}\begin{tabular}{ccccccccccccc}
\toprule
题号 & 1 & 2 & 3 & 4 & 5 & 6 & 7 & 8 & 9 & 10 & 11 & 12 \\
\midrule
答案 & D & A & B & B & B & C & D & C & B & A & B & D \\
\bottomrule
\end{tabular}\end{center}

\gkpart{填空题}
\begin{itemize}[leftmargin=2.6em]
  \item[\textbf{13.}] $\left(-\infty,\dfrac12\right)$。
  \item[\textbf{14.}] $90^\circ$。
  \item[\textbf{15.}] $\dfrac23$。
  \item[\textbf{16.}] ①④。
\end{itemize}

\gkpart{解答题}
\begin{itemize}[leftmargin=2.6em]
  \item[\textbf{17.}] (1) $p=\dfrac15$；(2) $\dfrac{243}{250}$
  \item[\textbf{18.}] (1) 最小正周期 $2\pi$，值域 $\left[-\dfrac{\sqrt2}2,\ \dfrac{\sqrt2}2\right]$；(2) $\dfrac7{25}$
  \item[\textbf{19.}] (1) $\dfrac{\sqrt3}4$；(2) $\dfrac{\sqrt5}5$
  \item[\textbf{20.}] (1) $a_n=\dfrac{2^n}{\lambda}$；(2) $n=6$
  \item[\textbf{21.}] (1) $4x^2-y^2=4$，即 $x^2-\dfrac{y^2}4=1$（$x\neq\pm1$）；(2) $(1,3)$
  \item[\textbf{22.}] (1) $f(n)=a^n$；(2) $a$ 的最小值为 $3$；(3) 前者大于后者，证明见解析
\end{itemize}
